\section{The cotangent direct images and the decisive period scalar}
\label{inv:sec:partA}

\subsection{Local logarithmic bundles and a duality calculation}

We use the marked threefold, with its arbitrary admissible
constant $c$.  Coefficients of coherent cohomology are complex.  In this
section $\omega_f=K_X\otimes f^*K_B^{-1}$,
$\mathcal L=\Omega_B^1(\log\Sigma)=\OO_B(1)$,
$\mathcal A=\Omega_X^1(\log D)$, and
$\mathcal E=\mathcal A/f^*\mathcal L$.  We write $\gamma,u$ also for
their images in $R^1f_*\OO_X$; this convention concerns cohomology
classes and does not identify them with differential-form representatives.

\begin{proposition}[The structure-sheaf input in the precise form used below]\label{inv:A1.0}\label{wb:thm:structure-direct-images-audit}\label{wb:lem:cusp-structure-sheaf-audit}
The direct images $R^qf_*\OO_X$, for $q=0,1,2$, are
$\OO$, $G=\OO\gamma\oplus\OO(-1)$, and $\OO(-1)$.
The inclusion $\OO\gamma\subset G$ is a subbundle.  At the cusp,
$(\gamma,u)$ is a holomorphic frame, and cup product identifies
$\bigwedge^2G$ with $R^2f_*\OO_X$.  In the cyclic covering coordinate
$s$, the action on $G$ has characters $(1,\zeta^{-a})$, where
$(m,a)=(3,2)$ or $(4,1)$ and $\zeta=e^{-2\pi i/m}$.
These assertions concern full scheme fibres, including $3S_1$ and $4S_2$.
At the cusp, the natural map $H^1(W;\mathbb C)\to H^1(W,\mathcal O_W)$ is an isomorphism. In degree two the comparison is surjective and is injective on $\ker[H^2(W;\mathbb C)\to H^2(\mathrm{dP}_6;\mathbb C)]$. The cup product $\bigwedge^2H^1(W,\mathcal O_W)\to H^2(W,\mathcal O_W)$ is an isomorphism. Every base-change map $(R^qf_*\mathcal O_X)\otimes k(b)\to H^q(X_b,\mathcal O_{X_b})$ is an isomorphism, including at the nonreduced fibres.
\end{proposition}
\Uses{the hypotheses, notation and definitions stated here; \citet[Lemma~9.12 and Proposition~9.13]{AC2026}; \cref{app:rederive:inv-A1.0}}


The structure-sheaf and cusp statements are due to \citet[Lemma~9.12 and Proposition~9.13]{AC2026}; we retain the independent root-line derivation. A complete re-derivation is given in Appendix~\ref{app:rederive:inv-A1.0}.


\begin{proposition}[Local freeness and the determinant line]\label{inv:A1.1}
Both $\mathcal A$ and $\mathcal E$ are locally free, of ranks three and
two.  The vertical vector $\xi=\partial/\partial\zeta_2$ is nowhere
zero as a section of $\mathcal E^\vee$.  There is an exact sequence
of holomorphic vector bundles
\begin{equation}\label{inv:eq:A-E-line}
0\longrightarrow N\longrightarrow\mathcal E
 \xrightarrow{\langle\cdot,\xi\rangle}\OO_X\longrightarrow0,
\qquad N=\det\mathcal E=\omega_f(-2S_1-3S_2).
\end{equation}
\end{proposition}
\Uses{the hypotheses, notation and definitions stated here; the explicit constructions and calculations in the proof}

\begin{proof}
At a smooth point of a reduced cusp branch the assertion is the usual
coordinate calculation.  At a double point write $t_c=xy$ with third
coordinate $z$; then
$\mathcal A=R\,d\log x\oplus R\,d\log y\oplus R\,dz$ and its
base subline is $R(d\log x+d\log y)$.  At either triple point write
$t_c=xyz$; then $\mathcal A$ has the three logarithmic coordinate
basis vectors and its base subline is their sum.  In each case that
subline is a direct summand.  These calculations apply on every one
of the three double curves, including their endpoints; no simplicity
of the global divisor is assumed.

At $S_j$ we use the actual free cyclic quotient
$q:J'_j\to N_j$, which is finite and etale on total spaces, and
$t_j=s^{m_j}$ on the base.  On $J'_j$ the logarithmic conormal is
$m_j\,ds/s$, and its coefficient is a unit over $\C$.  It is again a
direct summand of the logarithmic cotangent bundle.  The comparison
$q^*\mathcal A=\Omega_{J'_j}^1(\log(s=0))$ identifies the quotient
bundles and their exact sequence; faithful flat descent proves the
claim downstairs.  This is not a replacement of $m_jS_j$ by $S_j$.

On the marked torus family $\xi$ is the constant second vertical
vector.  The linear monodromies fix it, and translations preserve it.
On a cusp toric chart it is a nonzero constant vector in the relative
logarithmic tangent lattice of rank two.  A toric coordinate change
$x_i\mapsto c_i(t_c)t_c^{b_i}x_i$ changes $d\log x_i$ by a multiple
of $d\log t_c$, so this constant relative vector descends and is
nowhere zero in that bundle.  It can vanish as an ordinary tangent
vector at a toric stratum; this does not imply vanishing in the
logarithmic bundle.  Hence evaluation is surjective and its kernel
is the determinant line.  Taking determinants of
$0\to f^*\mathcal L\to\mathcal A\to\mathcal E\to0$ gives
\[
\det\mathcal E=K_X(D)\otimes f^*\mathcal L^{-1}
=\omega_f\bigl(D-f^*\Sigma\bigr)
=\omega_f(-2S_1-3S_2),
\]
which proves \eqref{inv:eq:A-E-line}.
\end{proof}

\begin{proposition}[The logarithmic direct images by vertical twisting and duality]\label{inv:A1.2}\label{wb:prop:log-relative-direct-images}
One has, as holomorphic vector bundles on $B$,
\[
\begin{array}{c|ccc}
q&0&1&2\\\hline
R^qf_*\mathcal E&\OO\oplus\OO(-1)&
 \OO^2\oplus\OO(-1)^2&\OO\oplus\OO(-1).
\end{array}
\]
They commute with base change to every full scheme fibre.
\end{proposition}
\Uses{the hypotheses, notation and definitions stated here; \cref{inv:eq:A-H-images}; \cref{inv:eq:A-E-line}; \cref{inv:eq:A-N-images}; \cref{inv:A1.0}}

\begin{proof}
Put $H=2S_1+3S_2$.  We first compute the line bundle $\OO_X(H)$,
rather than a tensor product of two separately descended lattices.
On the cover at $p_j$, its sections allow poles of order at most $m_j-1$
in $s$.  A fixed character has invariant exponent zero and gains no
pole, whereas the variable character $\zeta^{-a}$ has least allowed
invariant exponent $a-m$.  Thus the inclusion
$R^1f_*\OO\to R^1f_*\OO(H)$ is the identity on $\OO\gamma$ and
changes $s^ab_1$ to $s^{a-m}b_1=t_j^{-1}s^ab_1$ on its quotient.
It makes exactly one upper modification at each finite point.
The same computation applies to $R^2$.  No modification is made at
the cusp or a regular fibre.  The result is
\begin{equation}\label{inv:eq:A-H-images}
f_*\OO(H)=\OO,\qquad
R^1f_*\OO(H)=\OO\oplus\OO(1),\qquad
R^2f_*\OO(H)=\OO(1).
\end{equation}
Here the middle extension splits because $H^1(\OO(-1))=0$.
The local invariant modules are free.  Tensoring them with
$\C\{t_j\}/(t_j)$ gives invariants of
$s^{-(m_j-1)}\C\{s\}/(s^m)$, with the indicated surviving character
basis vectors.  The Reynolds idempotent is linear over the base,
so this calculation also proves the comparison with the nonreduced
fibre.  At the cusp the line bundle is trivial.  Cohomology and base
change, applied successively from the top degree, proves the claimed
local freeness and comparison for \eqref{inv:eq:A-H-images}.

For clarity the relative duality map we use is the cup product followed
by the analytic trace, giving a morphism of bundles
\[
R^qf_*(\omega_f\otimes\OO(-H))\longrightarrow
\bigl(R^{2-q}f_*\OO(H)\bigr)^\vee.
\]
Every fibre is a compact Gorenstein surface: it is the Cartier fibre
in a smooth threefold.  Serre duality on that scheme makes the map
an isomorphism on its residue field, also for $m_jS_j$.  The constant
dimensions just computed imply local freeness on both sides, so
Nakayama's lemma makes the bundle map an isomorphism.  Therefore
\begin{equation}\label{inv:eq:A-N-images}
(f_*N,R^1f_*N,R^2f_*N)
=(\OO(-1),\OO\oplus\OO(-1),\OO).
\end{equation}
This is the explicit reduction of analytic relative duality needed
here; the scheme fibre is retained throughout.

The connecting maps of \eqref{inv:eq:A-E-line} are zero on a regular
fibre: its relative cotangent is a trivial bundle and evaluation on
a constant nonzero vector splits.  Source and target direct images
in \eqref{inv:eq:A-N-images} and Statement \ref{inv:A1.0} are vector bundles,
so these holomorphic connecting maps, zero on a dense open set, are
zero everywhere.  We obtain extensions in each degree of the listed
$R^qN$ by $R^q\OO$.  Every resulting pair of line-bundle degrees is
zero or minus one; their extension groups are
$H^1(\OO(-1)),H^1(\OO),H^1(\OO(1))$, all zero.  These extensions
give exactly the three asserted bundles.  Applying the already proved
fibre comparisons to the short exact sequences proves base change
for $R^q\mathcal E$ as well.  This argument uses the vertical twist
$\OO(H)$ and duality; it does not use the workbench's filtration of
$f_*\mathcal E\otimes G$.
\end{proof}


\paragraph{Cited input.} \label{inv:A1.3}\label{wb:prop:absolute-cotangent-direct-images}\label{inv:eq:A-absolute-table}The source statement and its expanded derivation are retained in the proof compendium, Paper~I, label \nolinkurl{inv:A1.3}.\par


\subsection{An exact modular determinant and two numerical routes}

The scalar depends on the fourth root chosen in $t-1=s^4$.  We make
that ambiguity explicit.  We choose the local branch for which, on
the positive $t>1$ arc leading to the distinguished cusp,
$\tau=iy$ with $y>1$, and
\begin{equation}\label{inv:eq:A-s-branch}
\tau=i+\frac{i}{\pi\sqrt{E_4(i)}}s^2+O(s^4),\qquad
\operatorname{Im}\mu_s(0)>0.
\end{equation}
The nonvanishing proved below makes the second condition meaningful;
it fixes the remaining sign of $s$.  Equivalently it is the lower
cusp-regular zero branch of the Weierstrass function in Statement \ref{inv:A2.3}.
Replacing $s$ by $i^js$ multiplies the reported derivative by $i^{-j}$.
The workbench instead normalizes $ds/dz(z_2)=1$; its number is obtained
by the chain rule and is not numerically the same unspecified coordinate.

\begin{proposition}[The complete raw period determinant]\label{inv:A2.1}\label{wb:prop:log-ks-divisor}\label{wb:prop:period-derivative-determinant}
On the marked upper half-plane,
\begin{equation}\label{inv:eq:A-modular-K}
6\mu_z^2-\tau_z\beta_z
=\tau_z^2\frac{E_4(\tau)^2}{E_6(\tau)}.
\end{equation}
The right side extends holomorphically through the apparent poles.
Its exact orders are two at $z_1$ and zero at $z_2$; it has no other
zero away from their translates.  The logarithmic connecting map
$\kappa_0:f_*\mathcal E\to\mathcal L\otimes G$ has divisor
$\operatorname{divisor}(\det\kappa_0)=p_1+p_2$.
\end{proposition}
\Uses{the hypotheses, notation and definitions stated here; the explicit constructions and calculations in the proof}

\begin{proof}
Consider the holomorphic quadratic differential
$\mathscr Q=6(d\mu)^2-d\tau\,d\beta$.  Direct differentiation of
the two affine transformation laws in AC26 Definition 3.1 gives
\[
g_j^*\mathscr Q=\tau^{-2}\mathscr Q\quad(j=1,2).
\]
In particular $g_2^2$ fixes $\tau$ and sends
\[
\mu\longmapsto1-\tau-\mu,\qquad
\beta\longmapsto\beta+6\tau+12\mu-6,
\]
which leaves $\mathscr Q$ exactly unchanged.  The symbolic residuals
of all three identities are zero in \texttt{a\_calculations.py}.
The kernel of
$\Delta=\Z/3*\Z/4\to\Z/3*\Z/2=\operatorname{PSL}_2(\Z)$
is the normal subgroup generated by $g_2^2$.  The squared multiplier
of $\mathscr Q$ is the ordinary modular multiplier
$(c\tau+d)^{-2}$, so it is one on every kernel element.  Two generic
points of the marked cover with the same $\tau$ have the same $t$
and differ by an element of $\Delta$; its modular image fixes a
generic point, hence is the identity.  It follows that
$k=\mathscr Q/(d\tau)^2$ is a single-valued meromorphic function of
$\tau$, transforming as a modular form of weight two.  This kernel
argument is needed: mere equivariance of the separate period functions
would not prove descent.

At $z_1$, write a linearizing parameter $r$ with
$g_1r=\zeta r$, $\zeta=e^{-2\pi i/3}$.  If $\mathscr Q=K_rdr^2$
has first nonzero coefficient of order $n$, its transformation law
at the fixed point reads $\zeta^{n+2}=\rho^{-2}=\zeta$.
Thus $n\equiv2\pmod3$, and $K_r$ vanishes to order at least two.
The map $\tau$ is unramified there by AC26 Theorem 3.4(i).  At
$z_2$, $\tau-i$ has order two and $K_s$ is holomorphic; therefore
$k$ has at most a simple pole at $i$.  Away from these points
$\tau$ is unramified.  At the cusp $\mu$ and $\beta+\tau$ are
holomorphic in $q$, so $k=1+O(q)$.

The modular function $k E_6/E_4^2$ is therefore holomorphic at
$\rho$, at $i$, and at the cusp, as well as everywhere else.
The usual zero statements for $E_4,E_6$ (simple at $\rho,i$,
respectively) are the modular facts used in AC26 Section 3.1.
This function descends to the compact modular curve $X(1)=\Pone$,
so it is constant.  Its cusp value is one.  This proves
\eqref{inv:eq:A-modular-K} and the exact two local orders, without the
workbench's degree count for the coarse determinant line.

For completeness we translate to the actual coarse lattices.  Choose
cover eigenframes $e_0,e_1$ of $\mathcal E$ with characters
$1,\zeta^a$, and $b_0,b_1$ of $G$ with characters
$1,\zeta^{-a}$.  Their invariant bases are
$e_0,s^{m-a}e_1$ and $b_0,s^ab_1$.  The logarithmic base frame is
$\ell=dt_j/t_j=m\,ds/s$.  Changing determinants from these frames
to the marked ones gives
\[
\det\kappa_0=\epsilon(t_j)m^{-2}s^{m+2-2a}K_s,
\qquad \epsilon(0)\ne0.
\]
This equality follows by taking the two factors $s/m$ from replacing
$ds$ by $\ell$, the factor $s^{m-a}$ in the source, and the factor
$s^{-a}$ in the target.  Thus at $(m,a)=(3,2)$ the total order in
$s$ is $1+2=3$, and at $(4,1)$ it is $4+0=4$.
Both coarse zeros are simple.  At the cusp, the frame change has
no elementary divisor and $dZ/d\tau$ tends to
$\left(\begin{smallmatrix}0&1\\-1&0\end{smallmatrix}\right)$,
with determinant one.  Formula \eqref{inv:eq:A-modular-K} excludes any
further zero, proving the asserted divisor.
\end{proof}

\begin{proposition}[The evaluated scalar and the two local leading terms]\label{inv:A2.2}
In the coordinate \eqref{inv:eq:A-s-branch},
\begin{equation}\label{inv:eq:A-mu-value}
\mu_s(0)=i\,C_4,
\qquad
C_4=\frac{\sqrt{2/3}}{\pi E_4(i)^{1/4}}
=\frac{4\sqrt\pi}{3^{3/4}\Gamma(1/4)^2}>0.
\end{equation}
Numerically,
\[
C_4=0.2366093140807852246877051134504496062943051340858243372192243072\ldots.
\]
For $r^3=t$, with $r>0$ on the arc $0<t<1$ and $\tau$ travelling
from $\rho$ to $i$ along the unit circle, set
\[
\lambda=\frac{3e^{5\pi i/6}}{2\pi E_6(\rho)^{1/3}}.
\]
Then $\tau=\rho+\lambda r+O(r^2)$ and
\[
K_r=-\frac{3i\lambda}{2\pi}r^2+O(r^3).
\]
The leading coefficient has real and imaginary parts
\[
\begin{aligned}
\operatorname{Re}K_2&=0.08010221108753827140763650309934451243\ldots,\\
\operatorname{Im}K_2&=0.13874109940222333968402480721280640441\ldots.
\end{aligned}
\]
\end{proposition}
\Uses{the hypotheses, notation and definitions stated here; \cref{inv:eq:A-s-branch}; \cref{inv:eq:A-modular-K}; \cref{inv:A2.3}}

\begin{proof}
Ramanujan's differential relations give
$E_6'(i)=-\pi i E_4(i)^2$ and
$E_4'(\rho)=-(2\pi i/3)E_6(\rho)$.
Using $j=1728E_4^3/(E_4^3-E_6^2)$, we obtain
\[
j(\tau)-1728=-1728\pi^2E_4(i)(\tau-i)^2+O((\tau-i)^3).
\]
Since $j=1728(1+s^4)$, the chosen branch is precisely
\eqref{inv:eq:A-s-branch}.  Formula \eqref{inv:eq:A-modular-K} now gives
\[
K_s(0)=-\frac4{\pi^2\sqrt{E_4(i)}}=6\mu_s(0)^2,
\]
because $\tau_s(0)=0$.  This proves nonvanishing and the first
formula in \eqref{inv:eq:A-mu-value}; the sign is fixed by the stated
coordinate convention.  Similarly
$j(\tau)=-512i\pi^3E_6(\rho)(\tau-\rho)^3+cdots$.
This gives the stated $\lambda$, and substitution into
\eqref{inv:eq:A-modular-K} gives the coefficient $-3i\lambda/(2\pi)$.

There is also a geometric nonvanishing proof independent of this
modular determinant.  Using the section identification proved in
\ref{inv:A2.3}, set $x=s^2X,y=s^3Y,t=1+s^4$.  The smooth family becomes
\[
Y^2=X^3-3(1+s^4)^3X+2(1+s^4)^4s^2,
\qquad P=(0,-\sqrt2(1+s^4)^2s)
\]
on the chosen branch.  At $s=0$ its equation is
$Y^2=X^3-3X$ and the point is $(0,0)$.  The first variation of the
curve itself is zero.  In the local coordinate $Y$ at that point,
$dX/(2Y)=dY/(3X^2-3)$, so the Abel integral of $P$ has derivative
$\sqrt2/3$.  The marked period is $-i\Omega_R$, rather than the
positive real period $\Omega_R$: the positive imaginary branch
$\tau=iy$, $y>1$, has elliptic scale $F/s=i/\Omega_R$ at zero.
Its period has zero first derivative.  Division gives
$\mu_s(0)=i\sqrt2/(3\Omega_R)\ne0$.
Using $g_2(\Z+i\Z)=(4\pi^4/3)E_4(i)$ and
$g_2(\Omega_R(\Z+i\Z))=12$ gives
$\Omega_R=\pi E_4(i)^{1/4}/\sqrt3$, which identifies this value
with \eqref{inv:eq:A-mu-value}.  This also explains why choosing the
positive real period with a positive real $s$ instead would select
the reciprocal branch $\operatorname{Im}\tau<1$.

Here are the numerical methods actually run.  The first evaluates
$E_4(i)=1+240\sum_{n=1}^{100}\sigma_3(n)e^{-2\pi n}$ at 100
decimal working digits and inserts it in \eqref{inv:eq:A-mu-value}.
The second integrates
\[
\Omega_R=\frac{\sqrt2}{3^{1/4}}
\int_0^{\pi/2}(1-\tfrac12\sin^2\theta)^{-1/2}\,d\theta,
\qquad C_4=\frac{\sqrt2}{3\Omega_R}.
\]
This is the real period of $Y^2=X^3-3X$ for $dX/(2Y)$; substituting
the three roots $\sqrt3,0,-\sqrt3$ in the elementary elliptic-integral
change of variable proves the formula.  The beta-integral evaluation
of the displayed complete elliptic integral gives the gamma
expression in \eqref{inv:eq:A-mu-value}.  The Fourier and quadrature
answers agree with absolute discrepancy $1.79\cdot10^{-102}$.
As a third check, Newton iteration on the Weierstrass Fourier series
with $j(\tau)=1728(1+s^4)$, using the two roots adjacent to
$(1-\tau)/2$, gives their symmetric derivative.  At $s=10^{-15}$ it
is
\[
i\,0.23660931408078522468770511345048081721971218\ldots,
\]
with discrepancy $3.13\cdot10^{-32}$ and maximum residual below
$2.9\cdot10^{-101}$.  The order-three coefficient was independently
evaluated using the divisor-sum expansion of $E_6$ and its Jacobi
theta-constant formula, with agreement to the working precision.
The code, all output, and precision are supplied.  These are numerical
corroborations, not interval certificates; nonvanishing and the exact
orders were proved analytically above.  For example the elementary
bound $\sigma_3(n)\le n^4$ and $e^{-2\pi}<1/500$ bounds the omitted
Fourier tail after $N$ terms by
$250(N+1)^4 500^{-(N+1)}$; roundoff is separately controlled by
the working precision and the independent calculation.
\end{proof}

\begin{proposition}[Identification of the elliptic logarithm and the verified zero formula]\label{inv:A2.3}
For the Weierstrass equation
\[
y^2=x^3-3t^3(t-1)x+2t^4(t-1)^2,
\]
the section $[\mu]$ is a generator of its rational section group and
equals one of $P_*=(0,\sqrt2t^2(t-1))$ and $-P_*$.
In particular $\wp(\mu,\tau)=0$ on every regular fibre.
Along the branch in \eqref{inv:eq:A-s-branch}, with
$x_0=1-1/t$, the exact branch is
\begin{equation}\label{inv:eq:A-DI-branch}
\mu=\frac{1-\tau}{2}
 +\frac{i\sqrt6}{3\pi}
 \frac{x_0^{1/4}\,{}_3F_2(1/3,2/3,1;3/4,5/4;x_0)}
 {{}_2F_1(1/12,5/12;1;1-x_0)}.
\end{equation}
Consequently $\mu-(2-\tau)/3$ has exact order two at the
order-three point in $r^3=t$.
\end{proposition}
\Uses{the hypotheses, notation and definitions stated here; \cref{inv:eq:A-s-branch}; \citet{EichlerZagier1982}; \citet{DukeImamoglu2008}; \cref{inv:eq:A-mu-value}}

\begin{proof}
We include the comparison, since a matching $j$-map alone would not
identify a section.  Write $F$ for the elliptic scale, not the
threefold projection.  It is specified by
\[
F^2=-\frac{3t(t-1)E_4(\tau)}{\pi^2E_6(\tau)},\qquad
F/t\longrightarrow i\sqrt3/\pi
\]
at the distinguished cusp.  The local orders are two at $r=0$ and
one at $s=0$; hence the square root exists on the simply connected
marked cover.  Its transformation laws are
$F(g_1z)=-F(z)/\tau(z)$ and $F(g_2z)=F(z)/\tau(z)$.
These follow first after squaring by modularity, and the signs are
fixed by those local orders.  Direct substitution in the cubic
identifies its regular fibres with $\C/(\Z+\tau\Z)$ by
\[
x=F^2\wp(\nu,\tau),\qquad y=\tfrac12F^3\wp'(\nu,\tau).
\]
The transformation laws of $\mu$ make this a single-valued section
$P$.  Its coordinates are meromorphic at the two finite special
points and at the cusp, so are rational functions on $\Pone$.

Every rational section $Q$ admits a global holomorphic elliptic
logarithm $\nu$ on the marked cover.  Indeed the exponential map
from the trivial complex line bundle to this smooth elliptic family
is a covering map with locally uniformly discrete lattice; pullback
over a simply connected base lifts the section.  The apparent poles
in its projective Weierstrass coordinates at a finite point are
meromorphic poles of finite order: the coordinates of $Q$ are
rational in $t$, while $F$ is holomorphic with the finite orders
just computed.  Multiplying the three homogeneous coordinates by
a common power of the local uniformizer and cancelling their
common zero gives holomorphic coordinates not all zero at the
centre.  Their limiting point lies on the closed projective
Weierstrass family.  This is the valuative criterion for properness
over the local analytic discrete valuation ring; it is not an
assertion that an arbitrary punctured-disc holomorphic map to a
compact space has a removable singularity.
We can choose this logarithm bounded at the distinguished cusp.
To verify this last assertion, put $v=1/t$, $X=v^2x$, $Y=v^3y$.
The equation becomes
$Y^2=X^3-3(1-v)X+2(1-v)^2$.  A section cannot meet its nodal point
$(X,Y)=(1,0)$ at $v=0$: substituting $X=1+O(v),Y=O(v)$ leaves
nonzero coefficient one in the defining equation at order $v$.
On the smooth nodal locus the coordinate
$u=\exp(2\pi i\nu)$ is a nonzero finite coordinate, as follows
directly from
\[
X=1+\frac{12u}{(u-1)^2},\qquad
Y=\frac{12\sqrt3u(u+1)}{(u-1)^3}.
\]
The normally convergent Weierstrass Fourier series extends this
coordinate to a nearby smoothing.  The inverse function theorem
therefore writes the section as a holomorphic unit $u(v)$, also
near $u=1$, the point at infinity.  A logarithm of that unit fixes
the bounded lift.  Any two such lifts differ by an integer, since
an additional nonzero multiple of $\tau$ is unbounded.

For this lift, write the integral transformation errors as
\[
\nu(g_1z)=-\nu/\tau+\alpha_1+\beta_1(\tau-1)/\tau,
\quad
\nu(g_2z)=\nu/\tau+\alpha_2-\beta_2/\tau.
\]
Bounded cusp regularity and $g_1g_2=g_0^{-1}$ imply
$\alpha_2=-\beta_1$, $\beta_2=\alpha_1+\beta_1$.
Putting $n=-\alpha_1-2\beta_1$, $a=\alpha_1+\beta_1$, we find that
$\nu-(n\mu+a)$ satisfies the homogeneous law of AC26 Lemma 3.10
and is cusp-regular.  We recall a direct verification of the degree
used here.  The function
$F_h=E_4(\tau)^2E_6(\tau)^{1/2}/\Delta(\tau)$ on the marked cover
has orders two and one at $z_1,z_2$, and has a simple pole at the
cusp.  The square root exists because its zeros on this cover have
even order.  Modularity and these local orders fix the two signs:
$F_h(g_1z)=-F_h/\tau$, $F_h(g_2z)=F_h/\tau$.
For any holomorphic solution of the same homogeneous law, the
quotient by $F_h$ is an invariant meromorphic function of $t$.
A possible pole at the order-three point would have order at most
$2/3$ in $t$, and one at the order-four point at most $1/4$; its
integral order is therefore nonnegative.  At the cusp regularity
requires this quotient to vanish.  Thus the solution sheaf is
exactly $\OO_B(-p_0)$, with no nonzero global sections.
We obtain $\nu=n\mu+a$ and hence $Q=nP$.  The two distinct
finite values of $\mu$ exclude $nP=0$ for $n\ne0$ by cusp boundedness;
there is no torsion.  This proves that $P$ generates the section group.

We identify $P_*$ with a primitive generator by an intersection
calculation.  The rational elliptic surface and its fibre types
$\mathrm{IV}^*,\mathrm{III},I_1$ are those in Engel Section 2.
At $t=1$, the blowup chart $x=(t-1)X,y=(t-1)Y$ gives
$Y^2=(t-1)X^3-3t^3X+2t^4$; $P_*$ meets its exceptional
multiplicity-one component at $(X,Y)=(0,\sqrt2)$, away from the
identity component.  At $t=0$, two successive charts
$x=t^2X,y=t^2Y$ give
$Y^2=t^2X^3-3t(t-1)X+2(t-1)^2$; $P_*$ meets the smooth exceptional
component $Y=-\sqrt2$ at $X=0$.  These are the two nonidentity
simple components of the $E_6$ configuration.  The section is
disjoint from $O$ in every chart, including infinity.

Here the height calculation can be made directly, without assuming
primitivity.  Orthogonally project $Q-O$ in
$\operatorname{NS}(S)\otimes\Q$ off $O$, the fibre $F_S$, and the
nonidentity components.  The inverse intersection matrices are
$-C_{E_6}^{-1}$ and $-1/2$.  Their diagonal entries for the two
nonidentity simple components are $-4/3$ and $-1/2$.
Here we may number the nonidentity components so that
\[
C_{E_6}=\begin{pmatrix}
2&-1&0&0&0&0\\
-1&2&-1&0&-1&0\\
0&-1&2&-1&0&0\\
0&0&-1&2&0&0\\
0&-1&0&0&2&-1\\
0&0&0&0&-1&2
\end{pmatrix},\qquad
(C_{E_6}^{-1})_{44}=(C_{E_6}^{-1})_{66}=4/3.
\]
Since $K_S=-F_S$, adjunction gives $Q^2=O^2=-1$.  Expanding the
square of this projection gives
\[
h(Q)=2+2(Q\cdot O)-c_0(Q)-c_1(Q),\quad
c_0(Q)\in\{0,4/3\},\ c_1(Q)\in\{0,1/2\}.
\]
The projection is additive: the chord-and-tangent law makes
$Q+R-(Q+R)_{\rm section}-O$ linearly equivalent to a vertical
divisor, which the projection kills.  Its negative square therefore
satisfies $h(nQ)=n^2h(Q)$.  A section meets a single simple
component, since its base parameter has order one; consequently
the listed local corrections exhaust the possibilities.
For $Q\ne O$ we have $Q\cdot O\ge0$, so $h(Q)\ge1/6$.
For $P_*$ the explicit incidences give $h(P_*)=1/6$.
Writing $P_*=nP$ now gives $1/6=n^2h(P)\ge n^2/6$, and hence
$n=\pm1$.  Inverting the displayed $E_6$ Cartan matrix and the
chord-and-tangent formulas is included in the exact computation log.
This proves $\wp(\mu,\tau)=0$.

We checked the original theorem on page 399 of
\citet{EichlerZagier1982}, as well as equation (6) of
\citet{DukeImamoglu2008}.  With the normalized $\Delta$, their
coefficient is indeed $144\pi i\sqrt6$, their constant is
$\log(5+2\sqrt6)/(2\pi i)$, and the integrand is
$(\sigma-\tau)\Delta(\sigma)/E_6(\sigma)^{3/2}$ on the vertical
path with its continued branch.  Thus the displayed Eichler--Zagier integral formula is
correct with its indicated sign and period ambiguity.  Theorem 1
of Duke--Imamoglu gives \eqref{inv:eq:A-DI-branch}: we take one minus
their $z_0$, rather than $z_0-\tau$, because the latter is unbounded
at the distinguished cusp.  The chosen branch approaches
$(1-i)/2$ at $\tau=i$ and has lower imaginary cusp value.
These conditions identify the original bounded lift.  In particular
its linear coefficient is
$i\sqrt6/(3\pi\,{}_2F_1(1/12,5/12;1;1))$, agreeing with
\eqref{inv:eq:A-mu-value}.

Finally exact elliptic addition gives
\[
x(2P_*)=9t^2/8,\quad
y(2P_*)=\frac{t^2(32-5t)}{16\sqrt2},\quad
x(3P_*)=\frac{16}{81}(t-1)(5t-32).
\]
The last expression has nonzero limit $512/81$ at $t=0$.
If $\nu=\mu-(2-\tau)/3$, then $3\nu$ is the local logarithm of
$3P$ modulo its exact period.  Since $F$ has order two and
$\wp(3\nu)=(3\nu)^{-2}+O(\nu^2)$, the equality
$x(3P)=F^2\wp(3\nu)$ forces $\nu$ to have exact order two.
It also gives its leading coefficient
$\nu_2=\pm3F_2/(16\sqrt2)$, where $F_2=\lim_{r\to0}F/r^2$;
the sign is the explicitly stated elliptic-logarithm branch.
\end{proof}

\begin{proposition}[The global rank and the decisive scalar]\label{inv:A2.4}\label{wb:thm:h12-scalar}
The map
$\kappa_1:R^1f_*\mathcal E\to\mathcal L\otimes R^2f_*\OO=\OO$
is surjective, and $H^0(\kappa_1):\C^2\to\C$ has rank one.
Let $h=a\,d\zeta_1+d\zeta_2$ be the unique global section of
$f_*\mathcal E$ with $h(\xi)=1$.  In a genuine invariant target
frame normalized as below,
\[
\kappa_1(h\gamma)(p_2)=-\tfrac14\mu_s(0)=-\tfrac i4 C_4.
\]
The rank-one condition is equivalent to
$a\tau_z+\mu_z\not\equiv0$.  Under the same direct-image and
sheaf-surjectivity hypotheses a rank-zero global map would instead
give $h^{1,2}=2$, $h^{1,1}=3$; that alternative does not occur here.
\end{proposition}
\Uses{the hypotheses, notation and definitions stated here; \cref{inv:A2.1}; \cref{inv:A1.3}}

\begin{proof}
Let $d_{ij}(s)$ be the cover Kodaira--Spencer coefficients in the
eigenframes of Statement \ref{inv:A2.1} after removing $ds$.
The invariant coarse matrix and its cup-product row have the forms
\[
[\kappa_0]=\begin{pmatrix}A&t_jB\\C&D_0\end{pmatrix},\qquad
[\kappa_1]=(-C,A,-D_0,B),
\]
in the ordered basis
$e_0b_0,s^ae_0b_1,s^{m-a}e_1b_0,e_1b_1$.  Specifically
$A=sd_{00}/m$, $B=s^{1-a}d_{01}/m$,
$C=s^{1-a}d_{10}/m$, $D_0=s^{m+1-2a}d_{11}/m$.
These are holomorphic invariant functions; the row follows by
wedging the connecting class with $b_0,b_1$.
If the row had no unit at a finite point, every entry would vanish
modulo $t_j$ and $\det\kappa_0$ would vanish to order at least two,
contradicting \ref{inv:A2.1}.  Away from the finite points, including the cusp,
$\kappa_0$ is invertible, which gives a unit in the cup-product row.
Thus $\kappa_1$ is surjective.

The lift $h$ exists and is unique because $f_*\mathcal E$ is an
extension of $\OO$ by $\OO(-1)$.  Its invariance under $g_2$ gives
$a(g_2z)=\tau a+\mu$, hence $a(z_2)=1/2$.
In the quotient frame $(\gamma,u)$ of $G$ the connecting map is
\[
\kappa_0(h)=\frac{s}{4}
\bigl((6a\mu_s+\beta_s)\gamma+(a\tau_s+\mu_s)u\bigr)\ell.
\]
Choose the variable eigenframe $b_1$ by Reynolds averaging so that
$\gamma\wedge b_1=v(s)\gamma\wedge u$, $v(0)=1$.
Then $s\,v(s)(\gamma\wedge u)\ell$ is the actual invariant
target frame; the uncorrected expression with $v=1$ need not be
invariant away from zero.  Cup product with $\gamma$ gives at the
central point $-(a\tau_s+\mu_s)/4=-\mu_s/4$, proving rank one.

To check the stated equivalence, there is a second global section
\[
q_E=\sigma_1\otimes u+6\sigma_2\otimes\gamma
\quad\text{of }R^1f_*\mathcal E.
\]
It is the negative of the relative $(1,1)$ image of the invariant
flat class $\eta=u\wedge w+6\gamma\wedge\delta$.
It extends at the cusp because the relative logarithmic forms and
$\gamma,u$ extend there.  At a finite point it is an invariant
holomorphic tensor on the smooth cover, hence belongs to the four
explicit invariant generators above.  Its cup-product row is zero,
since the two terms are $6\mu_z$ and $-6\mu_z$.
It is independent of $h\gamma$ on every regular fibre (its
$\sigma_1\otimes u$ coefficient is nonzero), so these two sections
form a basis of the two-dimensional global section space.
Consequently the global map has rank one precisely when its value
on $h\gamma$, equivalently $a\tau_z+\mu_z$, is not identically
zero.  If its rank were zero, the long exact cohomology sequence of
$0\to\mathcal K\to\OO^2\oplus\OO(-1)^2\to\OO\to0$
would give $h^0(\mathcal K)=2$, $h^1(\mathcal K)=1$.
The residue calculation in \ref{inv:A1.3} and Leray would then give exactly
the two counterfactual Hodge numbers stated in the claim.
\end{proof}

\subsection{The Hodge table and its Leray generators}

\begin{proposition}[The complete Hodge table for every admissible parameter]\label{inv:A3.1}\label{wb:cor:exact-hodge-diamond}\label{wb:thm:parametric-hodge-diamond}\label{wb:prop:workbench-holomorphic-form-vanishing}
For every $\operatorname{Im}c<-M$, the Hodge numbers are
\[
\begin{array}{c|rrrr}
q\backslash p&0&1&2&3\\\hline
3&0&0&0&1\\
2&0&1&2&1\\
1&1&2&1&0\\
0&1&0&0&0.
\end{array}
\]
\end{proposition}
\Uses{the hypotheses, notation and definitions stated here; \cref{inv:eq:A-absolute-table}; \cref{inv:A1.0}}


The vanishing of holomorphic forms and the Hodge-diamond
pattern with one undetermined integer are due to
\citet[Proposition~9.20, p.~80; Theorem~9.1(vi), p.~65]{AC2026}.
The value of that integer is determined here by the nonzero direct-image
scalar. The full calculations of both the pattern and its value are retained.

\begin{proof}
Since $B$ is a curve, coherent Leray has no differential and gives
for $F=\Omega_X^p$ the exact sequences
$0\to H^1(B,R^{q-1}f_*F)\to H^q(X,F)\to H^0(B,R^qf_*F)\to0$.
Substitution of \eqref{inv:eq:A-absolute-table}, with
$h^0(\OO)=h^1(\OO(-2))=1$ and all cohomology of $\OO(-1)$ zero,
gives columns $p=1,2$ directly.  Statement \ref{inv:A1.0} gives column zero,
and Serre duality gives column three.  Thus the numbers are outputs
of the sheaf computation.  Replacing $\beta$ by $\beta+c$ leaves
every derivative, finite eigencharacter, invariant lattice and
logarithmic residue used above unchanged.  The toric gluing units
change, but their relative logarithmic derivatives still differ only
by base logarithmic forms.  Hence the same computations apply at
each admissible $c$, including every point arbitrarily close to the
strict boundary; no manifold at the excluded boundary is asserted.
\end{proof}

\begin{proposition}[Specified associated-graded generators and a global Dolbeault lift]\label{inv:A3.2}\label{wb:lem:bcpr-eta-extension}
The Leray associated-graded spaces for the four middle groups are
\[
\begin{array}{c|cc}
&E_2^{1,q-1}&E_2^{0,q}\\\hline
H^{1,1}&H^1(\OO(-2))=\C\alpha&\C\eta_{\rm rel}\\
H^{1,2}&0&\C\upsilon\\
H^{2,1}&\C\upsilon^\vee&0\\
H^{2,2}&\C\eta_{\rm rel}^\vee&\C\alpha^\vee.
\end{array}
\]
Here $\alpha$ is the Cech class of $dt/t$ on the two standard affine
base charts, $\eta_{\rm rel}=-q_E$ is the invariant relative class
from \ref{inv:A2.4}, and $\upsilon$ is the image of the trace-dual of
$1\in f_*\OO(H)$ under $R^2f_*N\to R^2f_*\mathcal E$ and the
isomorphism $R^2f_*\Omega_X^1=R^2f_*\mathcal E$.
The superscript $\vee$ denotes the bases dual under the relative
trace and the base Serre pairing.  There exists a Dolbeault class
$Q\in H^1(X,\Omega_X^1)$ whose restriction to every regular fibre
is $\eta=u\wedge w+6\gamma\wedge\delta$.
\end{proposition}
\Uses{the hypotheses, notation and definitions stated here; \cref{inv:A2.4}; \cref{inv:eq:A-E-line}; \cref{inv:A1.3}}

\begin{proof}
The isomorphism $f_*\Omega_X^1=\Omega_B^1$ in the residue sequence
is the actual pullback map, so $\alpha$ gives its stated Leray
class; it can be represented by a nonzero normalization of
$f^*\omega_{\rm FS}$ in Dolbeault cohomology.  The section
$\eta_{\rm rel}$ belongs to the canonical kernel identification
$R^1f_*\Omega_X^1=\ker\kappa_1$ and spans its global sections.
Leray is surjective onto this space, so any preimage defines $Q$
with the asserted fibre restriction.  Its ambiguity is exactly
$\C\alpha$; the statement does not claim a canonical global form.
The injection $R^2f_*N=\OO\to R^2f_*\mathcal E$ from
\eqref{inv:eq:A-E-line} spans its global sections, since the quotient
is $\OO(-1)$.  This defines $\upsilon$ with no choice of a negative
degree summand.  Duality used in \ref{inv:A1.3} gives the remaining entries
and their specified dual pairings.  Choosing a lift of
$\eta_{\rm rel}$ and of $\alpha^\vee$ gives bases of the whole
middle cohomology groups.  The associated-graded descriptions are
canonical with the named choices; the splitting of Leray vector
spaces is not asserted canonical.
\end{proof}

\subsection{Ordinary relative differentials and the correction at the cusp}

\begin{proposition}[All local modules needed to distinguish degrees one and two]\label{inv:A4.1}
Let $R=\C\{x,y,z\}$.  For $t=xy$,
\[
\Omega_{R/\C\{t\}}^1=R^3/R(y,x,0)
\simeq(x,y)\oplus R,
\]
and
\[
\Omega_{R/\C\{t\}}^2\simeq
(x,y)\oplus R/(x,y),\qquad
\Omega_{R/\C\{t\}}^3\simeq R/(x,y).
\]
For $t=xyz$,
$\Omega^1=R^3/R(yz,xz,xy)$ is torsion-free of rank two and is
not free at the triple point.  If $I=(xy,xz,yz)$, there is an exact
sequence
\[
0\to T\to\Omega^2\xrightarrow{\wedge dt} I\to0,
\qquad
T=\frac{R u\oplus R v}{(xu,\ y(u-v),\ zv)},
\quad IT=0,
\]
and $\Omega^3=R/I$.  Finally, for $t=s^m$ in
$R=\C\{s,y_1,y_2\}$,
\[
\begin{aligned}
\Omega^1&=R^2\oplus R/(s^{m-1})\,ds,\\
\Omega^2&=R\,dy_1\wedge dy_2
 \oplus\bigl(R/(s^{m-1})\bigr)^2,\\
\Omega^3&=R/(s^{m-1})\,ds\wedge dy_1\wedge dy_2.
\end{aligned}
\]
\end{proposition}
\Uses{the hypotheses, notation and definitions stated here; the explicit constructions and calculations in the proof}

\begin{proof}
The degree-one presentation is the quotient by $dt$; higher degrees
are its exterior powers.  At a double point the transverse quotient
$R^2/R(y,x)$ maps isomorphically to $(x,y)$ by
$(a,b)\mapsto xa-yb$ after an immaterial change of sign of one
basis vector.  More explicitly one may use
$(a,b)\mapsto xa-yb$ with relation $(y,x)$; its kernel is precisely
that relation, since $x,y$ are relatively prime.  In degree two,
with basis $A=dx\wedge dy,B=dx\wedge dz,C=dy\wedge dz$, the
relations are $xA=yA=0$, $yB+xC=0$.  This gives the displayed
torsion summand and the ideal $(x,y)$; wedging once more gives
$R/(x,y)$.

At the triple point the map from ordinary to logarithmic relative
one-forms sends $(a,b,c)$ to $(xa-zc,yb-zc)$.
Its kernel is precisely $R(yz,xz,xy)$: the two equations imply
$c=xyh$, $a=yzh$, $b=xzh$ by relative primality of the coordinates.
Thus it injects into $R^2$, proving torsion-freeness.  Its fibre at
the maximal ideal has dimension three although its rank is two,
so it is not free.  The same minimal-generator argument works at
the double locus.

In degree two, wedging with $dt$ is the row
$(xy,-xz,yz)$ in the basis $(A,B,C)$.
Its syzygies are generated freely by
$u=(z,y,0)$ and $v=(0,y,x)$.
To check generation, an equation $xyA-xzB+yzC=0$ implies
$x\mid C$, say $C=xC_1$; then $yA=z(B-yC_1)$ implies
$A=zA_1$, $B=y(A_1+C_1)$.  Thus the syzygy is $A_1u+C_1v$.
The three relations from $dt\wedge(dx,dy,dz)$ are
$-xu,y(u-v),zv$, proving the stated presentation of the kernel.
The ideal $I$ kills it by those relations.  Conversely the quotient
$I$ is torsion-free, so this kernel is exactly the torsion of
$\Omega^2$.  The presentation immediately gives $\Omega^3=R/I$.
All matrix identities have also been checked over $\Z[x,y,z]$ by
the supplied exact program.  For $t=s^m$ the sole degree-one
relation is $ms^{m-1}ds=0$; $m$ is a unit and taking exterior
powers gives the three formulas.
\end{proof}

\begin{proposition}[Refuted: The claimed degree-one torsion along the cusp]\label{inv:A4.2}
The assertion in AC26 Remark 9.21, PDF page 81, that
$\Omega_{X/B}^1$ has torsion along $W\cup S_1\cup S_2$ is false.
Its torsion is supported precisely on $S_1\cup S_2$; its degree-one
cusp modules are torsion-free but not locally free on the double
locus.  The plural statement in Remark 1.3, PDF page 5, must be
qualified by degree: $\Omega^2_{X/B}$ does have torsion on the cusp
double locus.  The workbench's degree-one correction is correct;
it must not be extended to degree two.
\end{proposition}
\Uses{the hypotheses, notation and definitions stated here; \cref{inv:A4.1}}

\begin{proof}
Both disputed sentences were checked on rendered pages of the
original PDF, so this is not a transcription correction.  The
double and triple charts in \ref{inv:A4.1} exhaust the singular strata of
$W$.  They prove absence of degree-one torsion at the cusp and
exhibit nonzero degree-two torsion there.  At a multiple fibre the
finite etale comparison
$q^*\Omega^1_{N_j/B}=\Omega^1_{J'_j/B}$ is an isomorphism,
and its last summand in \ref{inv:A4.1} is nonzero for $m=3,4$.
Faithful flatness descends both its existence and its support.
There are no other critical fibres and the ordinary relative
cotangent is locally free on their complement.  Thus the stated
support is exact.  The erroneous sentence is at transcription
line 4330 (the final paragraph of Remark 9.21); the workbench's
correct degree-one replacement occupies lines 980--1006 of
\texttt{hodge\_frolicher.tex}.  No numerical Hodge value in that
workbench is refuted by this calculation.
\end{proof}



\section{Fr\"olicher, Bott--Chern, Aeppli, and Hermitian metrics}
\label{inv:sec:partC}

\subsection{The differentials and their dependence on the Hodge number}

In this section $X$ is the standard marked $(3,4)$ threefold and
$h=h^{1,2}(X)=1$ by Part A.  We retain $h$ in the formal calculations
to specify exactly which conclusions depend on its evaluated value.
Every assertion involving $h$ is valid for every value
consistent with the Hodge table
\begin{equation}\label{inv:eq:C-hodge-input}
\begin{array}{c|rrrr}
q\backslash p&0&1&2&3\\\hline
3&0&0&0&1\\
2&0&h&h+1&1\\
1&1&h+1&h&0\\
0&1&0&0&0.
\end{array}
\end{equation}
Thus no conclusion below silently uses $h=1$.  The geometric inputs for
this part are the structure-sheaf calculation, the vanishing of
holomorphic forms, and the sphere topology, which give
\eqref{inv:eq:C-hodge-input}; Statements \ref{inv:A1.3} and \ref{inv:A3.1} determine $h$.
We use $d=\partial+\bar\partial$ on complex-valued
smooth forms and the complex orientation for integration.

\begin{proposition}[The complete Fr\"olicher differential in terms of $h$]\label{inv:C1.1}\label{wb:thm:frolicher-e2-candidate}
For a compact complex threefold with \eqref{inv:eq:C-hodge-input} and the
complex de Rham cohomology of $S^6$, the Fr\"olicher spectral sequence
satisfies $E_1\ne E_2=E_\infty$.  Its only nonzero $d_1$ ranks are
\[
\rank d_1^{0,1}=1,\quad \rank d_1^{1,1}=h,\quad
\rank d_1^{1,2}=h,\quad \rank d_1^{2,2}=1.
\]
In particular its two middle rows are exact.
\end{proposition}
\Uses{the hypotheses, notation and definitions stated here; \cref{inv:eq:C-hodge-input}}

\begin{proof}
Here $E_1^{p,q}=H^q(X,\Omega_X^p)$ and
$d_r:E_r^{p,q}\longrightarrow E_r^{p+r,q-r+1}$ are complex-linear.
For $r=2$, any source in the two middle rows has target in one of
$(2,0),(3,0),(3,1)$ or outside the table, except the source $(0,2)$,
which is zero.  Sources in row $q=3$ outside $(3,3)$ are zero, and
$(3,3)$ has no target.  For $r=3$ the only possible targets inside the
table have sources $(0,2)$ or $(0,3)$, both zero.  For $r>3$ there is
no arrow inside the table.  A zero term remains zero on subsequent
pages.  Hence $E_2=E_\infty$.

The filtration on $H^k_{\mathrm{dR}}(X;\C)$ has no nonzero associated
graded term for $1\le k\le5$.  Every middle term of $E_2$ therefore
vanishes.  Consequently
\[
0\longrightarrow\C\longrightarrow\C^{h+1}
\longrightarrow\C^h\longrightarrow0,
\qquad
0\longrightarrow\C^h\longrightarrow\C^{h+1}
\longrightarrow\C\longrightarrow0
\]
are exact, giving the stated ranks.  The first rank is one, so $E_1$
cannot equal $E_2$.
\end{proof}

\begin{proposition}[A basis on which all first differentials are specified]\label{inv:C1.2}
Choose a nonzero $g\in H^{0,1}_{\bar\partial}(X)$, a basis
$q_1,\ldots,q_h$ of a complement to $\C d_1g$ in
$H^{1,1}_{\bar\partial}(X)$, a basis $v_1,\ldots,v_h$ of
$H^{1,2}_{\bar\partial}(X)$, and a class $e$ complementary to the
image of $d_1:H^{1,2}\to H^{2,2}$.  With
$a=d_1g$, $c_j=d_1q_j$, $d_j=d_1v_j$, and $t=d_1e$, the first
differential is
\[
g\mapsto a\mapsto0,\quad q_j\mapsto c_j\mapsto0,
\qquad v_j\mapsto d_j\mapsto0,\quad e\mapsto t\mapsto0.
\]
These are bases in the corresponding bidegrees.  If $h=1$ and a
class $q$ restricts to the invariant fibre class $\eta$ described below,
we may take $q_1=q$, $v_1=gq$, and $e=q^2$, after scaling $g$ if
necessary.  Statements \ref{inv:A1.3} and \ref{inv:A3.2} prove these hypotheses for $X$,
so this geometric choice applies to the threefold in the stated problem.
\end{proposition}
\Uses{the hypotheses, notation and definitions stated here; \cref{inv:A1.3}; \cref{inv:A3.2}; \cref{inv:C1.1}; \cref{inv:C3.7}}

\begin{proof}
The first assertions follow from the injectivity and surjectivity
proved in \ref{inv:C1.1}.  For the last assertion write
$\eta=u\wedge w+6\gamma\wedge\delta$ in the marked integral
cohomology of a regular fibre.  Its complex-oriented square has
integral $-12$.  Its product with the $(0,1)$ image of $\gamma$ is
nonzero: the explicit period calculation in \ref{inv:C3.7} gives its nonzero
pairing with $d\zeta_2$.  Thus $gq$ and $q^2$ have nonzero fibre
restrictions.  The images under $d_1$ restrict to zero on the fibre,
because the fibre is a complex torus and its invariant representatives
are $d$-closed.  Therefore $q$ complements $d_1g$, $gq$ spans the
one-dimensional $H^{1,2}$, and $q^2$ complements $d_1H^{1,2}$.
The extension used in this last choice is the actual Leray lift in
Statement \ref{inv:A3.2}.  This argument uses that extension; it does not infer
it from a calculation on the smooth fibre alone.
\end{proof}

\subsection{A computation from the two definitions}

\begin{proposition}[Bott--Chern and Aeppli tables as functions of $h$]\label{inv:C2.1}\label{wb:thm:exact-bott-chern-aeppli}
With
\[
H^{p,q}_{BC}=
\frac{\ker\partial\cap\ker\bar\partial\subset A^{p,q}}
{\im\partial\bar\partial},\qquad
H^{p,q}_{A}=
\frac{\ker\partial\bar\partial\subset A^{p,q}}
{\im\partial+\im\bar\partial},
\]
the complete tables are
\begin{equation}\label{inv:eq:C-BC-A}
\begin{array}{c|rrrr}
q\backslash p&0&1&2&3\\\hline
3&0&0&1&1\\
2&0&h&2h&1\\
1&0&2&h&0\\
0&1&0&0&0
\end{array}_{BC}
\qquad
\begin{array}{c|rrrr}
q\backslash p&0&1&2&3\\\hline
3&0&0&0&1\\
2&0&h&2&0\\
1&1&2h&h&0\\
0&1&1&0&0.
\end{array}_{A}
\end{equation}
In particular, substitution of the value $h=1$ proved in \ref{inv:A3.1} gives
the workbench's numerical tables.  The totals are
$(1,0,2,2,2,2,1)$ and $(1,2,2,2,2,0,1)$, respectively.
\end{proposition}
\Uses{the hypotheses, notation and definitions stated here; \cref{inv:A3.1}; \citet[Theorem 3]{Stelzig2021}; \cref{inv:C1.1}; \citet[Proposition 6(1)]{Stelzig2021}; \cref{inv:eq:C-arrows}}

\begin{proof}
We apply the structure theorem for bounded double complexes over a
field, \citet[Theorem 3]{Stelzig2021}, to the actual algebraic double
complex $(A^{\bullet,\bullet}(X),\partial,\bar\partial)$.
The word bounded refers to the bidegrees, not to the dimensions of the
spaces of smooth forms, so that theorem applies here.  It gives a
direct-sum decomposition by isomorphisms of double complexes into
squares and zigzags.  An odd zigzag has one-dimensional total
cohomology.  All odd zigzags in intermediate total degree are therefore
excluded by $b_1=\cdots=b_5=0$.  An odd zigzag of length greater than
one would have at least two vertices of its cohomological total degree.
The extreme diagonals $p+q=0$ and $p+q=6$ contain just one bidegree
each, so the remaining odd zigzags are exactly one dot at $(0,0)$ and
one at $(3,3)$.

An even zigzag of length $2r$ contributes a nonzero $d_r$ to one of
the two Fr\"olicher spectral sequences.  \ref{inv:C1.1} and complex conjugation
exclude $r\ge2$ in both sequences.  Thus the remaining summands are
arrows.  The horizontal arrows, with their multiplicities, are
\begin{equation}\label{inv:eq:C-arrows}
\begin{array}{c|c|r}
\text{source}&\text{target}&\text{multiplicity}\\\hline
(0,1)&(1,1)&1\\
(1,1)&(2,1)&h\\
(1,2)&(2,2)&h\\
(2,2)&(3,2)&1.
\end{array}
\end{equation}
The vertical arrows are their reflections in $p=q$, with the same
multiplicities.  The correspondence with spectral-sequence ranks is
\citet[Proposition 6(1)]{Stelzig2021}; alternatively one computes
column cohomology on each horizontal arrow directly.

We now evaluate the displayed definitions, rather than use a formula
for the dimensions.  On an arrow, the source is killed by neither
the outgoing differential nor its total differential, whereas its
target is closed under both.  There is no composite
$\partial\bar\partial$.  Hence the target contributes one dimension
to Bott--Chern cohomology, and the source contributes one dimension
to Aeppli cohomology.  On a square, only the top right corner is closed
under both operators, and it is the image of
$\partial\bar\partial$ from the bottom left corner.  The other
three corners of the square are images of $\partial$ or
$\bar\partial$, while the bottom left corner is not in
$\ker\partial\bar\partial$.  Thus a square contributes to neither
cohomology.  A dot contributes once to each.  These computations
commute with direct sums, including the possibly infinite sum of
squares.  Counting the sources and targets in \eqref{inv:eq:C-arrows}
and its reflection gives \eqref{inv:eq:C-BC-A}.

The program \texttt{code/c\_double\_complex.py} performs a separate
kernel-and-image computation over $\Q$ and a source-and-target
count.  It includes an acyclic square with its signs, checks
$\partial\bar\partial+\bar\partial\partial=0$, and verifies the
two answers for $h=0,1,2,3,4,5$.  Its output is
\texttt{outputs/c\_double\_complex.txt}.  This is exact linear
algebra, not a floating-point check and not a determination of $h$.
\end{proof}

\begin{proposition}[Comparison with McHugh and Angella]\label{inv:C2.2}
The specialization of McHugh's Theorem 6.3, Proposition 6.4, equation
(7), and Section 7.1 to \eqref{inv:eq:C-hodge-input} agrees with \ref{inv:C2.1}.
The corresponding relations and Figure 5 in Angella's Section 4 also
agree.  No correction to these specific formulas is required.
\end{proposition}
\Uses{the hypotheses, notation and definitions stated here; \cref{inv:eq:C-hodge-input}; \cref{inv:C2.1}; \cref{inv:C1.1}; \citet[Theorem 6.3, Proposition 6.4, (7), pp. 20--25]{McHugh2018}; \cref{inv:eq:C-BC-A}; \citet[Section 4, pp. 13--16]{Angella2018}}

\begin{proof}
McHugh's quotient $K^{p,q}$ is zero in intermediate total degrees
on every arrow and square just used: a closed corner of an arrow is
the image of the incoming differential and a closed corner of a square
is in both images.  Thus $k^{p,q}=0$ for $0<p+q<6$.
Also $E_2^{3,1}=0$ by \ref{inv:C1.1} and $H^{1,1}(X;\R)=0$ because
$H^2(X;\R)=0$.  Reading the formulas in
\citet[Theorem 6.3, Proposition 6.4, (7), pp. 20--25]{McHugh2018}
therefore gives
\[
h_{BC}^{1,1}=2h^{0,1}=2,\qquad
h_{BC}^{1,2}=h^{1,2}=h,\qquad
h_{BC}^{2,2}=-1+(h+1)+h=2h.
\]
The remaining entries in his table are precisely the boundary entries
of \eqref{inv:eq:C-BC-A}.  In
\citet[Section 4, pp. 13--16]{Angella2018}, the relation
\[
h_{BC}^{2,2}=2h_{BC}^{2,1}-2h^{0,1}+2
\]
gives $2h$, and the boundary values reduce to those above because
$h^{2,0}=h^{0,2}=0$.  This comparison concerns exactly the formulas
used here; it does not certify every assertion in either article.
\end{proof}

\subsection{The algebra, line bundles, and metric consequences}

\begin{proposition}[The real Bott--Chern Chern-class map and a global $G$]\label{inv:C3.1}\label{wb:cor:picard-line-audit}
Let $G$ be a $\bar\partial$-closed smooth $(0,1)$ representative of
a nonzero $g\in H^1(X,\OO_X)$.  Set
\[
\mathsf A=[\partial G]_{BC},\qquad
\mathsf B=[\bar\partial\bar G]_{BC}.
\]
These classes form a complex basis of $H^{1,1}_{BC}(X)$.
With the exponential $z\mapsto e^{2\pi iz}$ and the Chern convention
$c_{1,BC}=[iF/(2\pi)]$, the map of real Lie groups is
\begin{equation}\label{inv:eq:C-Pic-map}
\Pic(X)=\{\exp(zg):z\in\C\}\longrightarrow H^{1,1}_{BC}(X),
\qquad \exp(zg)\longmapsto z\mathsf A+\bar z\mathsf B.
\end{equation}
It is injective and its image is the real subspace fixed by conjugation.

There is an operator-explicit global representative.  Normalize the
Fubini--Study form by $\int_{\Pone}\omega_{FS}=1$, put
$\theta=f^*\omega_{FS}$, choose any smooth Riemannian metric on $X$,
and let $\mathcal G_d$ be its Green operator for the real de Rham
Laplacian.  Then
\begin{equation}\label{inv:eq:C-G-Green}
\alpha=d^*\mathcal G_d\theta,\qquad G_0=\alpha^{0,1}
\end{equation}
satisfy $d\alpha=\theta$, $\bar\partial G_0=0$,
$[G_0]\ne0$, and $\exp([G_0])=f^*\OO_{\Pone}(1)$.
Restriction $H^1(X,\mathcal O_X)\to H^1(X_b,\mathcal O_{X_b})$ is injective for every scheme-theoretic fibre; on a smooth fibre its image is $\mathbb C\gamma^{0,1}$.
\end{proposition}
\Uses{the hypotheses, notation and definitions stated here; \citet[Theorem~10.5]{AC2026}; \cref{app:rederive:inv-C3.1}}


The Picard and fibre-restriction statement is due to \citet[Theorem~10.5]{AC2026}. A complete re-derivation is given in Appendix~\ref{app:rederive:inv-C3.1}.


\begin{proposition}[The factor $12$ with a fixed collar orientation]\label{inv:C3.2}\label{wb:lem:bcpr-degree-one-input}
Use the marked cohomology class $\gamma$ dual to
$\hat\gamma$, the analytic peripheral generators and the lift relation of
[AC26, Section 7.5], and put
$p=12l_0-4l_1-3l_2=\pm1$.  Let $g_\gamma$ be the unique coherent
Leray class whose fibre image is $\gamma^{0,1}$.  Then
\begin{equation}\label{inv:eq:C-factor12}
\exp(12g_\gamma)=f^*\OO_{\Pone}(-p).
\end{equation}
For the standard data $(l_0,l_1,l_2)=(0,1,-1)$, this is
$\exp(12g_\gamma)=f^*\OO(1)$.  For the opposite value $p=1$,
the exponent representing $f^*\OO(1)$ is $-12g_\gamma$.
\end{proposition}
\Uses{the hypotheses, notation and definitions stated here; \cref{inv:eq:C-collar-sum}; \cref{inv:C3.1}}

\begin{proof}
We include the comparison that fixes the sign.  A degree-one integral
class $k\gamma$ on a regular fibre extends over the cyclic filling of
order $m_j$ precisely when the affine generator $\hat g_j$ can be
assigned an integer value $n_j$ satisfying
$m_j n_j=k\gamma(v_j)=kl_j$.  Admissibility gives
$\gcd(l_j,m_j)=1$, so this is precisely $m_j\mid k$.
At the cusp $\gamma$ extends, as follows from its actual specialization
to $H^1(W;\Z)$ in [AC26, Theorem B.1 and Proposition 7.12].
The global monodromy invariants have rank one generated by $\gamma$.
It follows, by intersection of these three local images in the marked
local system, that
$H^0(B,R^1f_*\Z)=12\Z\gamma$.  This recovers the degree-one
part of [AC26, Proposition 7.26], without identifying the reduced and
non-reduced fibre: this is topological cohomology, unchanged by
nilpotents, whereas the coherent comparison uses the full fibre.

Let $U=\bigsqcup_iD_i$ be small discs around the special points and
let $V$ be their pair-of-pants complement with collars.  Lift
$12\gamma$ to integral classes $\Theta_i$ on $f^{-1}(D_i)$ and
$\Theta_V$ on $f^{-1}(V)$.  Existence over $V$ follows from
$H^2(V;\Z)=0$ and the low-degree Leray sequence; over each disc it
follows from the local images just computed.  Write
$a_j=\Theta_V(\hat\rho_j)$.  On the two finite analytic peripheral generators
the differences $\Theta_i-\Theta_V$ have periods
\[
n_1=4l_1-a_1,\qquad n_2=3l_2-a_2.
\]
At the cusp, the specified lift relation
$\widetilde a_0=t_\mu\hat\rho_2^{-1}\hat\rho_1^{-1}$,
with $\gamma(\mu)=l_0$, and the null-homotopy of
$\widetilde a_0$ in the cusp filling give
$n_0=-12l_0+a_1+a_2$.  Hence
\begin{equation}\label{inv:eq:C-collar-sum}
n_0+n_1+n_2=-p.
\end{equation}
These are the explicitly oriented boundary evaluations in [AC26,
Proposition 7.27, pp. 62--63], not an unspecified sign of a spectral
sequence differential.

Represent the lifts by real closed one-forms.  On a collar their
difference is the pullback of a real closed base form $b_i$ with
clockwise period $n_i$, modulo an exact form.  This follows from the
degree-one Leray sequence over the annulus, since the difference has
zero fibre restriction and the annulus has no second cohomology.
Subtract exact terms on smaller collars, extending their primitives
with cutoffs, so the equality of forms holds wherever a partition of
unity changes.  For a real base cutoff $\rho_i$ equal to one toward
$D_i$ and zero toward $V$, the resulting glued real one-form $\alpha$
on $X$ has
\[
d\alpha=f^*\vartheta,\qquad
\vartheta=\sum_i d\rho_i\wedge b_i.
\]
The form $\vartheta$ is a smooth real base two-form; thus
$\bar\partial\alpha^{0,1}=0$.  Its fibre restriction is the
$(0,1)$ image of $12\gamma$, so the coherent Leray edge isomorphism
$H^1(X,\OO_X)\simeq H^0(B,R^1f_*\OO_X)$ identifies
$[\alpha^{0,1}]=12g_\gamma$.

In a positively oriented polar coordinate $(r,\varphi)$ on $D_i$,
$b_i=-n_i\,d\varphi/(2\pi)$ modulo an exact form.
Because $\rho_i$ decreases from one to zero as $r$ increases,
\[
\int_{D_i\cap V}d\rho_i\wedge b_i
=-n_i\int_0^1d\rho_i=n_i.
\]
The exact part does not change this integral, since $d\rho_i$ is
supported inside the collar.  Therefore
$\int_B\vartheta=-p$ by \eqref{inv:eq:C-collar-sum}.  \ref{inv:C3.1} identifies
the Bott--Chern class of $\exp(12g_\gamma)$ with $[f^*\vartheta]$.
On $\Pone$ two real $(1,1)$ forms of the same integral differ by
$i\partial\bar\partial$ of a real function (solve the ordinary
scalar Laplace equation).  Hence this class is the Chern class of
$f^*\OO(-p)$, and injectivity in \ref{inv:C3.1} proves
\eqref{inv:eq:C-factor12}.  The program
\texttt{code/c\_sign\_and\_fibre.py} verifies the integer
intersection, \eqref{inv:eq:C-collar-sum}, and the collar integral exactly.
\end{proof}

\begin{proposition}[The complete Bott--Chern algebra, for every $h$]\label{inv:C3.3}\label{wb:thm:bcpr-square-zero-algebra}\label{wb:prop:bcpr-generators-are-bases}
The positive-total-degree ideal $N\subset H^{\bullet,\bullet}_{BC}(X)$
satisfies
\[
H^{\bullet,\bullet}_{BC}(X)=\C\cdot1\oplus N,\qquad N^2=0.
\]
Its nonzero bidegrees and dimensions are those of \ref{inv:C2.1}.  Thus the
square-zero conclusion is independent of the evaluated value of $h$.
For $h=1$, the nine classes specified in \eqref{wb:eq:bcpr-generators}, together with the unit, give bases of every nonzero Bott--Chern group.
\end{proposition}
\Uses{the hypotheses, notation and definitions stated here; \cref{inv:C2.1}; \cref{wb:eq:bcpr-generators}; \cref{inv:C1.2}; \cref{inv:C3.1}; \cref{inv:C1.1}}

\begin{proof}
Take global $\bar\partial$-closed representatives $G,Q_j,V_j,E$
of the classes $g,q_j,v_j,e$ of \ref{inv:C1.2}.  By the arrow computation in
\ref{inv:C2.1}, bases for the positive Bott--Chern groups consist of
\[
\partial G,\ \bar\partial\bar G;\qquad
\partial Q_j,\ \bar\partial\bar Q_j;\qquad
\partial V_j,\ \bar\partial\bar V_j;\qquad
\partial E,\ \bar\partial\bar E;\qquad \nu,
\]
where the pairs have bidegrees $(1,1)$, $(2,1),(1,2)$,
$(2,2)$, $(3,2),(2,3)$, respectively, and
$\int_X\nu=1$.  In particular the natural map
$H^{3,2}_{BC}\to H^{3,2}_{\bar\partial}$ is an isomorphism:
on its unique horizontal arrow the closed target maps nontrivially
to column cohomology.

Write $a=\partial G$, $b=\bar\partial\bar G$ and
$c_j=\partial Q_j$.  The graded Leibniz rule gives
\[
\partial\bar\partial(G\wedge\bar G)=-a\wedge b,
\quad
\partial\bar\partial(G\wedge\bar Q_j)
=-a\wedge\bar c_j,
\quad
\partial\bar\partial(\bar G\wedge Q_j)=b\wedge c_j.
\]
Consequently $\mathsf A\mathsf B=0$,
$\mathsf A\overline{\mathsf C_j}=0$, and
$\mathsf B\mathsf C_j=0$.

Let $f^*\OO(1)=\exp(\lambda g)$, where $\lambda\ne0$.
Its curvature may be chosen to be the pullback of a two-form on $B$,
whose square vanishes pointwise.  By \ref{inv:C3.1},
$\lambda^2\mathsf A^2+\bar\lambda^2\mathsf B^2=0$.
The Dolbeault image of $\mathsf B^2$ is zero.  Hence
$[a^2]_{\bar\partial}=0$.  But
$a^2=\partial(G\wedge a)$, and \ref{inv:C1.1} says that
$d_1:H^{1,2}\to H^{2,2}$ is injective.  Thus
$[G\wedge a]_{\bar\partial}=0$.  Choose a smooth $(1,1)$ form
$\Phi$ satisfying $G\wedge a=\bar\partial\Phi$.  Applying
$\partial$ proves $a^2=\partial\bar\partial\Phi$.
Conjugation gives $\mathsf B^2=0$ as well.

To treat $\mathsf A\mathsf C_j$, pair its Dolbeault image with the
generator $g\in H^{0,1}_{\bar\partial}$.  Since
$\bar\partial c_j=0$ and $G\wedge a=\bar\partial\Phi$,
\[
\int_X a\wedge c_j\wedge G
=-\int_X\bar\partial(\Phi\wedge c_j)=0.
\]
Dolbeault--Serre duality and the isomorphism
$H^{3,2}_{BC}\to H^{3,2}_{\bar\partial}$ give
$\mathsf A\mathsf C_j=0$.  Complex conjugation gives
$\mathsf B\overline{\mathsf C_j}=0$.

Every remaining product whose bidegree lies in $[0,3]^2$ lands in
$(3,3)$: it is a product of a $(1,1)$ class with a $(2,2)$ class,
or of a $(2,1)$ class with a $(1,2)$ class.  Each non-top generator
displayed above is $d$-exact, since its chosen primitive is
$\bar\partial$-closed, or is the conjugate of such a form.
The products therefore have integral zero.  Integration is an
isomorphism $H^{3,3}_{BC}\to\C$ by \ref{inv:C2.1} and the class of $\nu$,
so all these products vanish.  All other products exceed the complex
dimension in at least one bidegree.  This exhausts the possibilities.
\end{proof}

\begin{proposition}[Non-K\"ahler degrees and the precise degreewise statement]\label{inv:C3.4}
For $k=0,\ldots,6$ the non-K\"ahler degrees are
\begin{equation}\label{inv:eq:C-delta}
(\Delta^k)=(0,2,2+2h,4h,2+2h,2,0).
\end{equation}
If $h\ge1$, the two maps
$\bigoplus_{p+q=k}H^{p,q}_{BC}\to H^k_{\mathrm{dR}}(X;\C)
\to\bigoplus_{p+q=k}H^{p,q}_A$
are both isomorphisms exactly for $k=0,6$.
For the weaker assertion that a $d$-exact pure $d$-closed form of
degree $k$ is $\partial\bar\partial$-exact, the valid degrees are
$k=0,1,6$, and also $k=3,4$ if $h=0$.
\end{proposition}
\Uses{the hypotheses, notation and definitions stated here; \cref{inv:C2.1}}

\begin{proof}
The totals from \ref{inv:C2.1} are
\[
(h^k_{BC})=(1,0,2,2h,2h,2,1),\qquad
(h^k_A)=(1,2,2h,2h,2,0,1).
\]
Subtracting twice the Betti numbers gives \eqref{inv:eq:C-delta}.
On the two endpoint dots both comparison maps are isomorphisms;
all intermediate de Rham groups vanish.  The first assertion follows
by inspecting when both outside groups vanish.  The weaker assertion
is exactly injectivity of the first comparison map.  In intermediate
degree this is equivalent to vanishing of the Bott--Chern group alone,
which gives the stated list.  In degree zero there is no nonzero
exact form, and in degree six integration identifies Bott--Chern and
de Rham cohomology.  This distinguishes the two meanings of a
``$\partial\bar\partial$-lemma in a degree''; a vanishing Bott--Chern
group in degree one does not make $\Delta^1$ vanish.
\end{proof}

\begin{proposition}[Metric obstructions that do not require $h$]\label{inv:C3.5}
The threefold $X$ admits no balanced, strongly Gauduchon, K\"ahler,
locally conformally K\"ahler, or Hermitian-symplectic metric.
\end{proposition}
\Uses{the hypotheses, notation and definitions stated here; the explicit constructions and calculations in the proof}

\begin{proof}
Let $\theta=f^*\omega_{FS}$.  It is a nonzero semipositive exact
real $(1,1)$ form.  If $\omega$ were balanced, the real four-form
$\omega^2$ would be closed and
$\int_X\theta\wedge\omega^2>0$, because the integrand is positive
on the regular locus.  Exactness of $\theta$ and Stokes' theorem
give zero, a contradiction.

More generally, if $\omega$ were strongly Gauduchon, write
$\partial\omega^2=\bar\partial u$ with $u\in A^{3,1}$.
The real four-form $\Omega=\omega^2-u-\bar u$ is closed.
Bidegrees give $\theta\wedge\Omega=\theta\wedge\omega^2$,
so the same argument applies.  K\"ahler metrics are balanced.
If $d\omega=\ell\wedge\omega$ with a closed real Lee form
$\ell$, then $b_1=0$ gives $\ell=dv$, and
$d(e^{-v}\omega)=0$, excluding locally conformally K\"ahler metrics.
Finally any symplectic two-form on $X$ is exact because $b_2=0$.
Its cube has nonzero integral for its symplectic orientation, whereas
an exact cube has zero integral.  Thus $X$ is not symplectic and in
particular admits no Hermitian-symplectic metric.
\end{proof}


\paragraph{Cited input.} \label{inv:C3.6}\label{inv:eq:C-Ricci}The source statement and its expanded derivation are retained in the proof compendium, Paper~I, label \nolinkurl{inv:C3.6}.\par

\Uses{the hypotheses, notation and definitions stated here; \citet[Theorem~2]{Gauduchon1985}; \citet[Theorem~1.3]{STW2017}; \cref{app:rederive:inv-C3.6}}


This is the application of \citet[Theorem~2]{Gauduchon1985} and \citet[Theorem~1.3]{STW2017} with the computed Bott--Chern class. A complete re-derivation is given in Appendix~\ref{app:rederive:inv-C3.6}.


\begin{proposition}[The fibre obstruction to SKT and astheno-K\"ahler metrics]\label{inv:C3.7}
The threefold $X$ admits no pluriclosed (SKT) metric and no
astheno-K\"ahler metric.  The geometric input is the value $h=1$
and the global Dolbeault class $q$ in \ref{inv:A3.2}, whose restriction to a
regular fibre is $\eta=u\wedge w+6\gamma\wedge\delta$.
More generally, it is enough that the real image of
$H_A^{1,1}(X;\R)\to H_A^{1,1}(F;\R)$ be contained in
$\R\eta$ for one regular fibre $F$.
\end{proposition}
\Uses{the hypotheses, notation and definitions stated here; \cref{inv:A3.2}; \cref{inv:C1.2}; \cref{inv:C2.1}; \cref{inv:eq:C-eta-negative}}

\begin{proof}
We first compute the sign that matters.  Put
$y=\operatorname{Im}\tau>0$, $m=\operatorname{Im}\mu$, and
$\Delta=\operatorname{Im}\beta-6m^2/y<0$.
For the marked period rows
$\sigma_1=6\mu\gamma+\tau u+w$ and
$\sigma_2=\beta\gamma+\mu u+\delta$ we have
$\eta=u\wedge\sigma_1+6\gamma\wedge\sigma_2$.
This proves that the real class $\eta$ has type $(1,1)$.
The determinant of the four real period columns, with rows ordered
as $\operatorname{Re}\zeta_1,\operatorname{Im}\zeta_1,
\operatorname{Re}\zeta_2,\operatorname{Im}\zeta_2$, is
$y\Delta<0$.  Since $\eta^2=12\gamma\wedge u\wedge w\wedge\delta$,
\begin{equation}\label{inv:eq:C-eta-negative}
\int_{F,\mathrm{or}_{\C}}\eta^2=-12.
\end{equation}
Equivalently the invariant representative of $\eta$ has Hermitian
coefficient matrix
\[
\frac1{y\Delta}
\begin{pmatrix}\operatorname{Im}\beta&-6m\\-6m&6y\end{pmatrix},
\qquad \det=\frac6{y\Delta}<0,
\]
so it has one positive and one negative eigenvalue.  The exact
symbolic determinant and an exact rational example are checked in
\texttt{code/c\_sign\_and\_fibre.py} by both eigenvalues and
diagonal elimination.

The explicit invariant fibre representative of $\gamma^{0,1}$ is
\[
\Gamma=\frac{i}{2\Delta}
\left(d\bar\zeta_2-\frac myd\bar\zeta_1\right).
\]
Writing $Q_F$ for the invariant $(1,1)$ representative of $\eta$,
direct exterior multiplication gives
\[
Q_F\wedge\Gamma=
\frac1{4y\Delta}
d\zeta_1\wedge d\bar\zeta_1\wedge d\bar\zeta_2,
\quad
\int_FQ_F\wedge\Gamma\wedge d\zeta_2=-1.
\]
This also proves the nonvanishing used in \ref{inv:C1.2}.

When $h=1$, \ref{inv:C2.1} identifies $H_A^{1,1}(X)$ with the sum of the
one horizontal source and its vertical conjugate.  A representative
$Q$ of $q$ is nonzero in the horizontal source, since its fibre
restriction is nonzero whereas $d_1g$ restricts to zero.  Thus
$[Q]_A,[\bar Q]_A$ form a basis.  Their restrictions to $F$ both
equal the real class $\eta$.  The real image of the restriction map
is therefore $\R\eta$.

If $\omega$ were a pluriclosed positive real $(1,1)$ form, its
restriction to $F$ would be pluriclosed and positive.  Average this
restriction under translations of the compact complex torus $F$.
The average $\omega_{\mathrm{av}}$ is invariant, closed, and
strictly positive.  Translation preserves the Aeppli class: on a
torus the $\partial\bar\partial$-lemma identifies Aeppli
cohomology with the invariant $(p,q)$ forms, so the average is the
unique invariant representative of that class.  This identification
can also be checked Fourier mode by Fourier mode: each nonzero mode
is a Koszul square for two independent type components of its real
nonzero covector and contributes zero Aeppli cohomology.  Thus
$[\omega_{\mathrm{av}}]=r\eta$ for some real $r$.
But positivity gives $\int_F\omega_{\mathrm{av}}^2>0$, whereas
\eqref{inv:eq:C-eta-negative} gives $-12r^2\le0$, a contradiction.
In dimension three the astheno-K\"ahler equation
$\partial\bar\partial\omega^{n-2}=0$ is exactly the SKT equation.
\end{proof}




